\section*{Hoofdstuk \ref{ch:b2:climate-proxies} --- Klimaatarchieven en proxy's}

\begin{solution}{exo:b2:climate-proxies:1}
(a) $\delta = (27.5 \times (-30) + 1350 \times 0)/(1350 + 27.5) = -825/1377.5 = -0.60$ promille.
(b) De oceaan zou dan al het water bevatten: $-0.60$ promille, \num{0.60} promille lichter dan
nu. (c) $27.5/1377.5 = \qty{2.0}{\percent}$.
\end{solution}

\begin{solution}{exo:b2:climate-proxies:2}
(a) $T = 16.1 - 4.76\,(2.50 - (0 - 0.27)) = 16.1 - 4.76 \times 2.77 = 16.1 - 13.19 =
\qty{2.9}{\celsius}$. (b) $T = 16.1 - 4.76\,(2.50 - 0.73) = 16.1 - 8.43 = \qty{7.7}{\celsius}$: dezelfde
schaal geeft een temperatuur die bijna vijf graden hoger is als het water zwaarder was; daarom
moet de samenstelling van het water bekend zijn. (c) De schaal is $2.71 - 2.50 = 0.21$ promille zwaarder,
dus haar water was $4.76 \times 0.21 = \qty{1.0}{\celsius}$ kouder, op \qty{1.9}{\celsius}.
\end{solution}

\begin{solution}{exo:b2:climate-proxies:3}
(a) $0.4\,\mathrm{e}^{0.9} = \qty{0.98}{mmol/mol}$ en $0.4\,\mathrm{e}^{1.8} =
\qty{2.42}{mmol/mol}$. (b) De factor is $\mathrm{e}^{0.9} = 2.46$; per graad $\mathrm{e}^{0.09}
= 1.094$, een groei van \qty{9.4}{\percent} per graad. (c) $T = \ln(\text{Mg/Ca}/0.4)/0.09$:
$\ln 5/0.09 = \qty{17.9}{\celsius}$ en $\ln 2.5/0.09 = \qty{10.2}{\celsius}$.
\end{solution}

\begin{solution}{exo:b2:climate-proxies:4}
(a) $2668 - 2213 + 1 = 456$ jaar. (b) Eén boom kan plaatselijk een ring missen, of een ring tonen
die niet te lezen is; het afstemmen van de patronen van brede en smalle ringen over veel bomen brengt de
ontbrekende ringen aan het licht en legt het jaar van elke ring vast, en de overeenstemming tussen de bomen meet de
betrouwbaarheid van de datering. (c) Aan de boomgrens beperkt één factor, de koude van het groeiseizoen,
alle bomen, zodat ze hetzelfde patroon van goede en slechte jaren delen; in een mild, vochtig
laagland overheerst geen enkele factor, verschillen de patronen van de bomen, en kan \omterm{def:b2:climate-proxies:dendro}{kruisdatering} mislukken.
\end{solution}

\begin{solution}{exo:b2:climate-proxies:5}
(a) $125 \times 40/(3682.2 - 125) = 5000/3557.2 = 1.41$ promille; $134 \times 40/3548.2 =
5360/3548.2 = 1.51$ promille. (b) Met $H - h = 3552.2$: $3900/3552.2 = 1.10$, $4550/3552.2 = 1.28$
en $6500/3552.2 = 1.83$ promille. (c) $\delta_{\mathrm{ice}} = -1.0 \times 3552.2/130 = -27.3$ pro%
mille, zwaarder dan elke huidige ijskap (Groenland $-30$ tot $-35$, de Noord-Amerikaanse ijskap
$-40$, Antarctisch ijs tot $-50$). De massabalans met realistisch ijs en de poriewater%
meting komen niet overeen; het verschil van ongeveer een halve promille is een echte onzekerheid van de methode.
\end{solution}

\begin{solution}{exo:b2:climate-proxies:6}
(a) $0.372\,\mathrm{e}^{0.0936T} = 0.4\,\mathrm{e}^{0.09T}$ geeft $\mathrm{e}^{0.0036T} =
0.4/0.372 = 1.0753$, dus $T = \ln 1.0753/0.0036 = \qty{20.2}{\celsius}$. (b) Bij
\qty{2}{\celsius}: $0.4\,\mathrm{e}^{0.18} = 0.479$ en $0.372\,\mathrm{e}^{0.1872} = 0.449$, dus de
aangepaste verhouding is $0.449/0.479 = 0.937$, \qty{6.3}{\percent} lager. (c) $T = \ln(0.449/0.4)/0.09 =
\ln 1.1225/0.09 = \qty{1.3}{\celsius}$: de voorspelde wet leest de schaal ongeveer
\qty{0.7}{\celsius} te koud. In koud water doet de keuze van de kalibratie ertoe op het niveau van een
graad.
\end{solution}

\begin{solution}{exo:b2:climate-proxies:7}
(a) Van het Holoceen tot GI-1: \qty{112.19}{m} ijs voor 2989 jaar: $11\,219/2989 =
\qty{3.75}{cm}$ per jaar; van GI-16 tot GI-17: \qty{17.89}{m} voor 1160 jaar: $1789/1160 =
\qty{1.54}{cm}$ per jaar. (b) Het ijs
vloeit: terwijl elke laag wordt begraven, spreidt ze zijwaarts uit en wordt ze dunner, en de diepere, oudere lagen zijn
langer uitgedund. (c) Wanneer de lagen te dun worden, een centimeter of minder, kunnen de instrumenten
het ene jaar niet meer van het volgende onderscheiden, en moet het \omterm{def:b2:climate-proxies:age-model}{ouderdomsmodel} in plaats daarvan ankerpunten of stromings%
modellen gebruiken.
\end{solution}

\begin{solution}{exo:b2:climate-proxies:8}
(a) Het deel van het interval is $(1550 - 1492.45)/(1604.64 - 1492.45) = 57.55/112.19 = 0.513$,
dus de ouderdom is $11\,703 + 0.513 \times 2989 = 11\,703 + 1533 = \num{13236}$ jaar b2k. (b) 1950 ligt 50
jaar vóór 2000: \num{13186} jaar vóór 1950. (c) $50 + 0.513 \times (93 - 50) = 72$ jaar.
\end{solution}

\begin{solution}{exo:b2:climate-proxies:9}
(a) Elke onzekere laag draagt een half jaar bij: $2 \times 93 = 186$ en $2 \times 1287 = 2574$
onzekere lagen. (b) $\frac12\sqrt{186} = 6.8$ jaar en $\frac12\sqrt{2574} = 25.4$ jaar.
(c) $93/6.8 = 14$ en $1287/25.4 = 51$. De verhouding van $\frac12 n$ tot $\frac12\sqrt{n}$ is
$\sqrt{n}$: ze groeit als de vierkantswortel van het aantal onzekere lagen, dus de keuze tussen
de twee regels doet er met de diepte steeds meer toe.
\end{solution}

\begin{solution}{exo:b2:climate-proxies:10}
(a) Op de \omterm{prop:b2:isotope-geochemistry:mwl}{meteorische waterlijn} verandert $\delta$D acht keer zoveel als $\delta^{18}$O, dus de
helling is $8 \times 0.8 = 6.4$ promille $\delta$D per graad. De daling is $441.7 - 396.6 = 45.1$
promille, dus een afkoeling van $45.1/6.4 = \qty{7.0}{\celsius}$. (b) De glaciale oceaan was ongeveer
$8 \times 1.0 = 8$ promille zwaarder in $\delta$D; sneeuw die uit de huidige oceaan ontstond, zou evenveel
lichter zijn geweest, $-441.7 - 8 = -449.7$ promille. De daling wordt $53.1$ promille en de afkoeling
$53.1/6.4 = \qty{8.3}{\celsius}$. (c) $9.39 - 0.09 = \qty{9.30}{\celsius}$, dus $45.1/9.30 = 4.85$
promille per graad. De gepubliceerde reeks is ook gecorrigeerd voor de veranderende hoogte van de ijs%
kap, en haar temperatuurbetrekking komt uit modellen die de waterisotopen simuleren, niet uit
de ruimtelijke helling alleen.
\end{solution}

\begin{solution}{exo:b2:climate-proxies:11}
(a) De logaritmen van de verhoudingen zijn $-0.528$, $-0.223$, $0.191$, $0.501$, $0.904$, $1.227$; hun
gemiddelde is $0.345$ en de gemiddelde temperatuur 14. Met $S_{TT} = \sum (T - 14)^2 = 280$ en $S_{T\ell} =
\sum (T - 14)(\ell - 0.345) = 24.93$ is $A = 24.93/280 = \qty{0.0890}{\per\celsius}$ en $\ln B =
0.345 - 0.0890 \times 14 = -0.901$, dus $B = \qty{0.406}{mmol/mol}$. (b) De residuen zijn
$+0.017$, $-0.034$, $+0.023$, $-0.023$, $+0.025$ en $-0.009$, enkele procenten van de verhouding, zonder
trend; de aanpassing ligt dicht bij $0.4\,\mathrm{e}^{0.09T}$. (c) $T = (\ln 1.30 + 0.901)/0.0890 = (0.262 +
0.901)/0.0890 = \qty{13.1}{\celsius}$. (d) $T = (\ln 0.45 + 0.901)/0.0890 = (-0.799 +
0.901)/0.0890 = \qty{1.2}{\celsius}$. Deze waarde ligt onder het koudste monster van de kalibratie,
\qty{4}{\celsius}: de wet wordt geëxtrapoleerd, en bij lage temperaturen wijkt de Mg/Ca van schalen bekend%
erwijs af van de exponentiële wet onder invloed van het carbonaation.
\end{solution}

\begin{solution}{exo:b2:climate-proxies:12}
(a) Per eenheid oceaanoppervlak wordt het water van de huidige oceaan, $H$ bij 0, verdeeld tussen de voorraden
$h_i$ bij $\delta_i$ en de resterende oceaan $H - h$, met $h = \sum h_i$: $0 = \sum h_i\delta_i +
(H - h)\,\Delta\delta$. Na de eerste stap is $\Delta\delta = 60 \times 30/(3682.2 - 60) =
1800/3622.2 = 0.50$ promille; na de tweede $(1800 + 70 \times 45)/(3682.2 - 130) =
4950/3552.2 = 1.39$ promille. (b) $-4950/130 = -38.1$ promille. (c) De regel is $1.39/130 = 0.0107$
promille per meter, en geeft $0.50/0.0107 = \qty{47}{m}$ voor de eerste stap in plaats van
\qty{60}{m}: het eerste ijs was zwaarder dan het gemiddelde, dus elke meter ervan verschoof de oceaan
minder. Een lineaire regel die op één toestand is gekalibreerd, leest de andere verkeerd wanneer de gemiddelde samenstelling van het
ijs verandert.
\end{solution}

\begin{solution}{pb:b2:climate-proxies:1}
\textbf{1.} $361.84 \times 10^6\ \unit{km^2} \times \qty{0.130}{km} = \qty{4.70e7}{km^3}$, ongeveer 47
miljoen kubieke kilometer.
\textbf{2.} $130 \times 40/(3682.2 - 130) = 5200/3552.2 = 1.46$ promille.
\textbf{3.} $5000/3557.2 = 1.41$ en $5360/3548.2 = 1.51$ promille.
\textbf{4.} $4550/3552.2 = 1.28$ en $6500/3552.2 = 1.83$ promille. De samenstelling van het ijs
verschuift het antwoord met \num{0.55} promille, het zeeniveau met \num{0.10}: het ijs weegt zwaarder.
\textbf{5.} $\delta_{\mathrm{ice}} = -1.0 \times 3552.2/130 = -27.3$ promille.
\textbf{6.} $40/3682.2 = 0.0109$ promille per meter, wat afgerond de \num{0.011} van de gedateerde
zeeniveaus geeft.
\textbf{7.} Bij constante $\delta^{18}\mathrm{O}_{\mathrm{sw}}$ is $\mathrm{d}T = -4.76\,\mathrm{d}\delta^{18}
\mathrm{O}_{\mathrm{b}}$, dus $\mathrm{d}T = -1$ geeft $\mathrm{d}\delta^{18}\mathrm{O}_{\mathrm{b}} =
1/4.76 = 0.21$ promille.
\textbf{8.} $\delta^{18}\mathrm{O}_{\mathrm{b}} = (16.1 - T)/4.76 + \delta^{18}\mathrm{O}_{\mathrm{sw}}
- 0.27 = 14.1/4.76 - 0.27 = 2.96 - 0.27 = 2.69$ promille.
\textbf{9.} $17.1/4.76 + 1.0 - 0.27 = 3.59 + 0.73 = 4.32$ promille. De verandering is $4.32 - 2.69 =
1.63$ promille: \num{1.00} door het water en $3.0 \times 0.21 = 0.63$ door de afkoeling van
\qty{3}{\celsius}.
\textbf{10.} $T = \ln(2.0/0.4)/0.09 = \ln 5/0.09 = \qty{17.9}{\celsius}$.
\textbf{11.} $0.18/0.1 = \qty{1.8}{\celsius}$ afkoeling.
\textbf{12.} Mg/Ca is een verhouding van twee elementen in het calciet, bepaald door de temperatuur van
de verkalking; ze betrekt de zuurstofisotopen niet, en negeert dus de samenstelling van
het water die door het ijs verandert.
\textbf{13.} $4.92 - 3.23 = 1.69$ promille.
\textbf{14.} $\Delta = h \times 40/(H - h)$ oplossen naar $h$: $h = H\Delta/(\Delta + 40) = 3682.2
\times 1.69/41.69 = \qty{149}{m}$, meer dan de 125 tot \qty{134}{m} van de zeeniveaugegevens:
een deel van de verandering moet afkoeling zijn.
\textbf{15.} $1.69 - 1.46 = 0.23$ promille, en $4.76 \times 0.23 = \qty{1.1}{\celsius}$.
\textbf{16.} $1.69 - 1.00 = 0.69$ promille, en $4.76 \times 0.69 = \qty{3.3}{\celsius}$ (3.28).
\textbf{17.} $0.23/0.25 = \qty{0.9}{\celsius}$ en $0.69/0.25 = \qty{2.8}{\celsius}$.
\textbf{18.} Het Noord-Atlantisch Diepwater koelde af van 2--\qty{4}{\celsius} tot \qty{-1.1}{\celsius},
met 3.1 tot \qty{5.1}{\celsius}; het Antarctisch Bodemwater van 0--\qty{1}{\celsius} tot
\qty{-1.9}{\celsius}, met 1.9 tot \qty{2.9}{\celsius}; de Mg/Ca van de noordwestelijke Atlantische Oceaan geeft
\qty{1.8}{\celsius}. Alle liggen tussen ongeveer 2 en \qty{5}{\celsius}, rond de 2.8 tot
\qty{3.3}{\celsius} van de poriewatercorrectie, en boven de \qty{1}{\celsius} van de massa%
balans met ijs van $-40$ promille.
\textbf{19.} $-396.6 - (-441.7) = 45.1$ promille.
\textbf{20.} $8 \times 0.8 = 6.4$ promille per graad, dus $45.1/6.4 = \qty{7.0}{\celsius}$.
\textbf{21.} $9.39 - 0.09 = \qty{9.30}{\celsius}$, en $45.1/9.30 = 4.85$ promille per graad.
\textbf{22.} Dome C, \qty{9.3}{\celsius}, meer dan het aardoppervlak, 5 tot \qty{7}{\celsius}, meer
dan de diepe oceaan, ongeveer \qty{3}{\celsius}.
\textbf{23.} Het diepe water is vandaag al koud, 0 tot \qty{4}{\celsius}, en zeewater kan niet
veel onder zijn vriespunt afkoelen, dat het glaciale bodemwater bij \qty{-1.9}{\celsius} bijna
bereikte: de diepe oceaan had maar enkele graden te verliezen, terwijl de lucht boven Antarctica
zo'n bodem niet heeft.
\textbf{24.} Met de poriewatercorrectie was de diepe oceaan tijdens het Laatste Glaciale Maximum ongeveer
\textbf{\qty{3.3}{\celsius} kouder} dan nu, dicht bij de 2 tot \qty{5}{\celsius} van de
onafhankelijke schattingen voor het diepe water.
\end{solution}
